\section{Discrete Cleanup, Dueling DDQN, Exploration, And Diagnostics}

\subsection{Scope}
This note documents the repo-wide reinforcement-learning cleanup performed for the unified MPC notebooks. The main goals are:
\begin{itemize}
  \item clean the standard discrete horizon baseline into a proper Double DQN,
  \item keep the dueling horizon family as a separate upgraded discrete path,
  \item add correct NoisyNet support for discrete exploration,
  \item expose configurable loss and exploration choices in the continuous-control agents,
  \item normalize mismatch features consistently, and
  \item save enough diagnostics in \texttt{input\_data.pkl} to regenerate plots later.
\end{itemize}

\subsection{Standard Horizon Baseline: Clean DDQN}
The standard horizon baseline now uses one online Q-network and one target Q-network. The target is the usual Double DQN construction:
\[
y_t = r_t + \gamma Q_{\theta^-}\!\left(s_{t+1}, \arg\max_a Q_{\theta}(s_{t+1}, a)\right)(1-d_t).
\]
This removes the older dormant twin-discrete-Q behavior and keeps evaluation aligned with the online network:
\[
a_t^{\mathrm{eval}} = \arg\max_a Q_{\theta}(s_t, a).
\]

The standard polymer and distillation horizon notebooks remain the baseline family. Their defaults stay conservative:
\begin{itemize}
  \item \texttt{exploration\_mode = "epsilon"}
  \item \texttt{loss\_type = "huber"}
  \item the existing unified horizon runner semantics are preserved
\end{itemize}

\subsection{Dueling DDQN Horizon Family}
The dueling path keeps the standard value--advantage decomposition:
\[
Q(s,a) = V(s) + A(s,a) - \frac{1}{|\mathcal{A}|}\sum_{a'} A(s,a').
\]
The target is still Double DQN:
\[
y_t = r_t + \gamma Q_{\theta^-}\!\left(s_{t+1}, \arg\max_a Q_{\theta}(s_{t+1}, a)\right)(1-d_t).
\]

Both polymer and distillation now have unified dueling horizon notebooks. The dueling defaults are research-facing:
\begin{itemize}
  \item \texttt{exploration\_mode = "noisy"}
  \item \texttt{loss\_type = "huber"}
  \item the same deep \texttt{[512,512,512,512,512]} backbone used for the standard horizon family
\end{itemize}

\subsection{NoisyNet Exploration}
Discrete exploration now supports epsilon-greedy and NoisyNet. The shared noisy layer uses factorized Gaussian noise:
\[
W = \mu_W + \sigma_W \odot \epsilon_W,\qquad
b = \mu_b + \sigma_b \odot \epsilon_b.
\]
Training and action collection use sampled noise when \texttt{exploration\_mode = "noisy"}, while evaluation disables noise and remains deterministic. This is implemented once and reused by both the standard DDQN and dueling DDQN stacks.

\subsection{Continuous-Agent Exploration And Loss Options}
The continuous stacks were extended without changing runner semantics:
\begin{itemize}
  \item TD3 now supports \texttt{loss\_type = "huber" | "mse"} and \texttt{exploration\_mode = "gaussian" | "param\_noise"}.
  \item Parameter-noise exploration perturbs only the data-collection actor and is resampled on a configurable cadence.
  \item SAC now supports \texttt{loss\_type = "huber" | "mse"} and saves richer training diagnostics.
  \item SAC evaluation stays deterministic through the mean action, while training remains stochastic.
\end{itemize}

\subsection{Mismatch-State Normalization}
Mismatch mode now scales innovation and tracking error by an explicit per-output supervisory range vector instead of using fixed raw clipping. The default scale comes from the existing setpoint-deviation bounds:
\[
\texttt{mismatch\_scale}_i = \max\!\left(|y_{\mathrm{sp,min},i}|,\; |y_{\mathrm{sp,max},i}|,\; \varepsilon\right).
\]
After normalization, an optional clip is applied, with default \texttt{mismatch\_clip = 3.0}. This preserves the original standard state path while making mismatch features comparable across systems and outputs.

\subsection{Diagnostics And Saved Bundles}
The saved run bundles now store richer metadata and diagnostics, including:
\begin{itemize}
  \item method family, algorithm, notebook source, run mode, state mode, config snapshot, and seed,
  \item discrete diagnostics such as loss, exploration, epsilon, TD error, max-Q, chosen-action Q, noisy sigma, and dueling value/advantage summaries,
  \item continuous diagnostics such as actor/critic losses, critic-Q traces, exploration magnitude, parameter-noise scale, action saturation, and SAC entropy / log-prob summaries.
\end{itemize}

The plotting layer was extended in a backward-compatible way so these saved bundles can drive:
\begin{itemize}
  \item richer discrete and continuous diagnostics dashboards,
  \item multi-run reward summaries with mean/std overlays,
  \item simple summary tables exported from saved \texttt{input\_data.pkl} bundles.
\end{itemize}

\subsection{Repo-Level Outcome}
The repo now has:
\begin{itemize}
  \item a clean DDQN baseline in \texttt{DQN/},
  \item an upgraded dueling DDQN stack in \texttt{DuelingDQN/},
  \item a polymer dueling unified horizon notebook,
  \item a distillation dueling unified horizon notebook,
  \item shared NoisyNet layers,
  \item configurable exploration/loss settings in DDQN, dueling DDQN, TD3, and SAC,
  \item richer result bundles and non-breaking plot extensions.
\end{itemize}

The existing unified runner structure and notebook families remain intact; this pass upgrades the learning components and diagnostics without redesigning the supervisory MPC workflow.

\subsection{N-Step Credit Assignment Extension}
The same unified stacks now support plain uncorrected $n$-step returns as an additive extension, with \texttt{n\_step = 1} kept as the canonical baseline. A shared accumulator forms endpoint transitions of the form
\[
\left(s_t,\; a_t,\; R_t^{(n)},\; s_{t+n},\; d_t^{(n)},\; \gamma_t^{(n)}\right),
\]
where
\[
R_t^{(n)} = \sum_{k=0}^{n_t-1} \gamma^k r_{t+k},
\qquad
\gamma_t^{(n)} =
\begin{cases}
0, & \text{if a terminal transition occurs before the bootstrap endpoint},\\
\gamma^{n_t}, & \text{otherwise.}
\end{cases}
\]
Here $n_t$ is the actual realized horizon, which may be smaller than the configured \texttt{n\_step} near the end of an episode.

This extension is threaded through:
\begin{itemize}
  \item the standard DDQN baseline in \texttt{DQN/},
  \item the dueling DDQN path in \texttt{DuelingDQN/},
  \item TD3 in \texttt{TD3Agent/},
  \item SAC in \texttt{SACAgent/},
  \item the active polymer and distillation unified notebooks for horizon, matrices, weights, and residual control.
\end{itemize}

The endpoint targets preserve each algorithm's original bootstrap rule while replacing the one-step reward/discount pair:
\begin{itemize}
  \item DDQN / dueling DDQN:
  \[
  y_t = R_t^{(n)} + \gamma_t^{(n)} Q_{\theta^-}\!\left(s_{t+n}, \arg\max_a Q_{\theta}(s_{t+n}, a)\right).
  \]
  \item TD3:
  \[
  y_t = R_t^{(n)} + \gamma_t^{(n)} \min_{i \in \{1,2\}} Q_{\phi_i^-}\!\left(s_{t+n}, \tilde a_{t+n}\right),
  \]
  with the same target-policy smoothing used in the one-step implementation.
  \item SAC:
  \[
  y_t = R_t^{(n)} + \gamma_t^{(n)} \left(\min_{i \in \{1,2\}} Q_{\phi_i^-}(s_{t+n}, a_{t+n}) - \alpha \log \pi(a_{t+n}\mid s_{t+n})\right).
  \]
\end{itemize}

The saved bundles now include \texttt{n\_step} metadata and simple decomposition traces such as mean aggregated reward, mean effective discount, mean bootstrap value, mean realized step count, and the truncated-fraction trace. This keeps old bundles backward-compatible while allowing $n$-step analyses to be reproduced later from \texttt{input\_data.pkl} alone.

\subsection{Advanced Multistep Stage 1}
The repo now extends the plain $n$-step path with a first-stage multistep research layer while preserving the original one-step defaults:
\begin{itemize}
  \item \texttt{multistep\_mode = "one\_step"} remains the default everywhere,
  \item \texttt{n\_step = 1} remains the canonical baseline,
  \item combined notebooks are intentionally left out of this first implementation pass.
\end{itemize}

Stage 1 adds three shared components:
\begin{itemize}
  \item contiguous local sequence sampling that respects episode boundaries,
  \item truncated forward-view $\lambda$-return target builders,
  \item exact discrete Retrace($\lambda$) support for epsilon-greedy DDQN and dueling DDQN only.
\end{itemize}

\paragraph{Sequence sampling.}
The new sequence samplers operate on local replay windows without crossing terminal states or episode resets. They return sequence masks and realized lengths so truncated finite-horizon targets remain explicit and do not silently leak across boundaries.

\paragraph{Lambda returns.}
For DDQN, dueling DDQN, TD3, and SAC, the finite-horizon $\lambda$ target is constructed from replay sequences with explicit truncation. At a high level,
\[
G_t^{\lambda} = r_t + \gamma \bigl((1-\lambda)\hat V_{t+1} + \lambda G_{t+1}^{\lambda}\bigr),
\]
with the recursion truncated by either the sampled sequence end or an episode termination.

\paragraph{Discrete Retrace.}
For epsilon-greedy DDQN and dueling DDQN, the new Retrace target uses clipped importance ratios
\[
\rho_t = \frac{\pi(a_t \mid s_t)}{\mu(a_t \mid s_t)},
\qquad
c_t = \lambda \min(1,\rho_t),
\]
and a backward correction term built over the sampled replay segment. This mode is intentionally restricted to epsilon-greedy exploration because the behavior probability is exactly known there. If a user selects Retrace with NoisyNet exploration, the code now fails fast with a clear configuration error rather than silently using an invalid approximation.

\paragraph{SAC-n.}
SAC now exposes \texttt{multistep\_mode = "sac\_n"} as the explicit research-facing name for its endpoint multistep soft backup, while preserving the same soft-target structure at the bootstrap endpoint:
\[
y_t = R_t^{(n)} + \gamma_t^{(n)}\left(\min_i Q_{\phi_i^-}(s_{t+n}, a_{t+n}) - \alpha \log \pi(a_{t+n}\mid s_{t+n})\right).
\]
The generic \texttt{"n\_step"} mode remains available as the simpler endpoint multistep label, and \texttt{"lambda"} is now also supported for critic-target construction.

\paragraph{Diagnostics and saved bundles.}
The result bundles and plotting layer now tolerate and consume the new optional multistep traces:
\begin{itemize}
  \item \texttt{multistep\_mode}, \texttt{lambda\_value},
  \item \texttt{lambda\_return\_mean\_trace},
  \item \texttt{offpolicy\_rho\_mean\_trace},
  \item \texttt{offpolicy\_c\_mean\_trace},
  \item \texttt{behavior\_logprob\_mean\_trace},
  \item \texttt{target\_logprob\_mean\_trace},
  \item \texttt{retrace\_c\_clip\_fraction\_trace}.
\end{itemize}
Old bundles remain valid because these fields are optional.

\paragraph{Deferred work.}
V-trace is intentionally deferred in this stage. The current implementation prepares the shared target and sequence infrastructure needed for broader off-policy sequence methods later, but does not expose V-trace as a selectable notebook mode yet.
